\section{Question, conventions, and scope}
% Canonical source-aligned PageRank/RPPR reference formulation.
%
% This fragment is intentionally heading-free at the section level so that the
% active paper and every standalone note can input it. It is a manuscript
% reference convention, not an implementation-wide residual decision.
% Primary source: Fountoulakis and Martinez-Rubio, arXiv:2602.21138v2,
% PDF p. 3, Section 3 and equation (RPPR).

\subsection{Common source-aligned PageRank and RPPR model}
\label{subsec:shared-source-aligned-problem}

All documents in this manuscript workspace use the following reference
language. It follows the plain italic notation of the comparison paper:
\(x,s,A,D,Q,I\) denote column vectors or matrices as determined by their
displayed domains. This is the scoped exception to the project's usual
bold-vector and bold-matrix typography. A note may impose a stronger
assumption after this block, but it must not redefine these objects.

Let \(G=(V,E)\) be a finite, simple, undirected, unweighted graph without
isolated vertices, let \(V=[n]:=\{1,\ldots,n\}\), and let
\(A\in\{0,1\}^{n\times n}\) be its symmetric adjacency matrix. For
\(i\in V\), write
\[
    \mathcal N(i):=\{j\in V:\{i,j\}\in E\},
    \qquad
    d_i:=|\mathcal N(i)|,
    \qquad
    D:=\diag(d_1,\ldots,d_n).
\]
For \(S\subseteq V\), define
\[
    \mathcal N(S):=\bigcup_{i\in S}\mathcal N(i),
    \qquad
    \partial S:=\mathcal N(S)\setminus S,
    \qquad
    \operatorname{Ext}(S):=V\setminus(S\cup\partial S),
\]
\begin{equation}
    \vol(S):=\sum_{i\in S}d_i.
    \label{eq:shared-volume}
\end{equation}
The support of a vector is
\(\supp(x):=\{i\in[n]:x_i\neq0\}\). Matrix inequalities use the Loewner
order, and all unqualified vector inequalities are coordinatewise.

The seed is a column distribution
\begin{equation}
    s\in\R_+^n,
    \qquad
    \inner{\one}{s}=1.
    \label{eq:shared-seed}
\end{equation}
The single-seed specialization is always written \(s=e_v\), where \(v\in V\)
is the seed vertex. Thus \(s\) is never used as a vertex label. For
\(\alpha\in(0,1]\), define
\begin{equation}
\begin{aligned}
    \mathcal L&:=I-D^{-1/2}AD^{-1/2},\\
    Q&:=\alpha I+\frac{1-\alpha}{2}\mathcal L
      =\frac{1+\alpha}{2}I
       -\frac{1-\alpha}{2}D^{-1/2}AD^{-1/2},\\
    b&:=\alpha D^{-1/2}s.
\end{aligned}
    \label{eq:shared-pagerank-matrices}
\end{equation}
Since \(0\preceq\mathcal L\preceq2I\),
\(\alpha I\preceq Q\preceq I\). The shared smooth PageRank quadratic is
\begin{equation}
    f(x):=\frac12\inner{x}{Qx}-\inner{b}{x},
    \qquad
    \nabla f(x)=Qx-b.
    \label{eq:shared-pagerank-objective}
\end{equation}
Its unique unregularized minimizer and associated lazy personalized PageRank
vector are
\begin{equation}
    x^0:=Q^{-1}b,
    \qquad
    \pi:=D^{1/2}x^0.
    \label{eq:shared-ppr-solution}
\end{equation}

For a regularization parameter \(\rho>0\), reserve \(g_\rho\) for the
nonsmooth weighted \(\ell_1\) term and define
\begin{equation}
    g_\rho(x):=\alpha\rho\norm{D^{1/2}x}_1,
    \qquad
    F_\rho(x):=f(x)+g_\rho(x),
    \qquad
    x^\star(\rho):=\argmin_{x\in\R^n}F_\rho(x).
    \label{eq:shared-rppr-objective}
\end{equation}
When \(\rho\) is fixed, \(x^\star\) abbreviates \(x^\star(\rho)\), never the
unregularized point \(x^0\). Write
\[
    S^\star(\rho):=\supp(x^\star(\rho)),
    \qquad
    I^\star(\rho):=[n]\setminus S^\star(\rho).
\]
The source records \(x^\star(\rho)\geq0\),
\(\vol(S^\star(\rho))\leq1/\rho\), and the coordinatewise conditions
\begin{equation}
    \nabla_i f(x^\star(\rho))
    \in
    \begin{cases}
        \{-\alpha\rho\sqrt{d_i}\}, & x_i^\star(\rho)>0,\\
        [-\alpha\rho\sqrt{d_i},0], & x_i^\star(\rho)=0.
    \end{cases}
    \label{eq:shared-rppr-kkt}
\end{equation}

The common computational unit is one repeated adjacency-list scan: scanning
coordinate \(i\) costs \(d_i\), and scanning a set \(S\) costs \(\vol(S)\).
Iteration count and wall-clock time do not replace this degree-weighted work
measure.

Finally, accuracy symbols remain distinct. The APPR threshold is
\(\epsappr\); \(\epsobj\) denotes an objective-gap target;
\(\epspg\) denotes the source experiment's proximal fixed-point tolerance; and
\(\epsppr\) may denote a document-scoped degree-normalized PPR error target.
No equality or conversion among these quantities is assumed. In particular,
this shared model does not resolve the repository-wide residual decision.

\subsection{Exact target and archived notation}
We take the point seed $\bm{s}=\bm{e}_v$ and objective tolerance
$\epsobj>0$. The target is an output $\widehat{\bm{x}}$ with
\[
 \bm{0}\le\widehat{\bm{x}}\le\bm{x}_\rho^*,\qquad
 F_\rho(\widehat{\bm{x}})-F_\rho(\bm{x}_\rho^*)\le\epsobj,
 \qquad\vol(\supp\widehat{\bm{x}})\le1/\rho.
\]
The cost includes scalar operations and comparisons, state maintenance,
first and repeated adjacency entries, degree queries, boundary processing,
and output. There is no supplied support, free preprocessing, or hidden
inverse oracle. The graph is accessed through its adjacency lists.

The archived proof uses ordinary italic letters for vectors and matrices.
Its $Q,D,b,x_r^*$ are respectively the shared $\bm{Q},\bm{D},\bm{b},\bm{x}_r^*$.
Its $w$ is $\bm{w}=\bm{D}^{1/2}\bm{1}$. Its objective tolerance means $\epsobj$.
Crucially, its stage source $s=b-Q\bar x$ is denoted here by
$\bm{h}=\bm{b}-\bm{Q}\overline{\bm{x}}$; it is not the shared seed
distribution $\bm{s}$. These are explicit note-scoped notation translations,
not changes of mathematical convention.

\paragraph{Proved here, with the complete proof archived.}
The result below refers to Theorem~1 of the original independently developed
proof, whose unmodified sources and complete audit records accompany this note.
The remaining sections summarize its dependency chain, rather than replacing
the archived detailed proofs.
\begin{theorem}\label{thm:archived-op2}
In the shared exact-real algebraic word model, for $0<\rho<1/d_v$, the
deterministic algorithm returns the three guarantees above in
$\widetilde O(1/(\rho\sqrt\alpha))$ fully charged work. Hidden factors are
polylogarithmic in inverse parameters, inverse objective accuracy, and exposed
or emitted words. For unrestricted $\rho>0$, add a constant to the bound.
\end{theorem}
If $\rho d_v\ge1$, zero is exact. If $\alpha=1$, only the seed coordinate
is needed. The rest of this note concerns the nontrivial case $0<\alpha<1$.
The rational-input bounded-integer realization is a separate extension;
ordinary unchecked floating-point arithmetic is not covered by that theorem.
